\section{Ergebnisse und Diskussion}
\label{sec:results}

\subsection{Inhaltliche Stichpunkte}
\begin{itemize}
  \item \textbf{Baseline gegenüber Full Context}
  \begin{itemize}
    \item Im Notebook~07 steigt Macro-F1 von 0{,}911805 (Baseline) auf
          0{,}953685 (Model~2); Accuracy steigt von 0{,}930625 auf
          0{,}962204 und ROC-AUC von 0{,}986353 auf 0{,}994932.
    \item Log Loss sinkt von 0{,}182478 auf 0{,}098996.
    \item Diese Werte sprechen im ausgewerteten Testlauf für bessere
          Klassifikation mit Full Context. Eine statistische Signifikanz
          ist damit ohne zusätzliche Tests nicht belegt.
  \end{itemize}
  \item \textbf{Ergebnisse nach Klasse}
  \begin{itemize}
    \item Backdoor: Recall 0{,}758200 $\rightarrow$ 0{,}880133;
          F1 0{,}810302 $\rightarrow$ 0{,}899922.
    \item none: Recall 0{,}963883 $\rightarrow$ 0{,}975301;
          F1 0{,}944577 $\rightarrow$ 0{,}967988.
    \item syn-flood: Recall 0{,}983818 $\rightarrow$ 0{,}993130;
          F1 0{,}980535 $\rightarrow$ 0{,}993145.
    \item Backdoor weist in der Baseline den niedrigsten Recall auf und
          zugleich den größten absoluten Recall-Anstieg im Full-Context-Modell.
  \end{itemize}
  \item \textbf{Zusätzliche Modellvariante}
  \begin{itemize}
    \item Die transformierte Baseline, bei der eigene Merkmale doppelt so
          stark wie aggregierte Merkmale gewichtet werden (Verhältnis 2:1),
          erreicht Accuracy 0{,}706217,
          Macro-F1 0{,}691625 und ROC-AUC 0{,}870686.
    \item Als separate Ablation/Transformation einordnen; nicht mit der
          Baseline ohne Kontext oder mit Model~1 gleichsetzen.
  \end{itemize}
  \item \textbf{Szenarien und Datenqualität}
  \begin{itemize}
    \item Notebook~07 erzeugt Szenariovergleiche für Distraction Attack~1/2,
          Distributed SYN Flood, Worm Backdoor und Cluster Shuffling
          SYN Flood.
    \item Notebook~01b zeigt beim hybriden Trainingsdatensatz gegenüber
          Originaltraining höhere Macro-F1- und ROC-AUC-Werte auf dem dort
          verwendeten Original-Testdatensatz; diese Datengenerierungs-
          Evaluation von der P2P-Evaluation getrennt diskutieren.
  \end{itemize}
\end{itemize}

\subsection{Noch zu erklären}
\begin{itemize}
  \item Die Falsch-Positiv-Rate ist nicht direkt in der Notebook-Tabelle
        enthalten. Aus Konfusionsmatrizen je Klasse berechnen und die
        Definition für das Mehrklassenproblem angeben, bevor Robustheit
        gegenüber Sensitivität bewertet wird.
  \item Für die Forschungsfrage den Unterschied zwischen beobachtetem
        Leistungsanstieg und statistischer Signifikanz herausarbeiten;
        gegebenenfalls gepaarte Tests oder Wiederholungen ergänzen.
  \item Model~1 in Notebook~06 wird mit 24\,000 Testfällen je Knoten
        (192\,000 insgesamt) berichtet, während der primäre Vergleich in
        Notebook~07 30\,000 je Knoten (240\,000 insgesamt) verwendet.
        Diese Auswertungen nicht direkt vermischen; kanonische Testmenge
        und Modellvariante festlegen.
  \item Die Szenarioplot-Ergebnisse pro Modell und Angriff erklären; keine
        Aussagen über einzelne Szenarien aus der aggregierten Gesamtmetrik
        ableiten.
  \item Für den Hybrid-Datensatz erklären, ob die Verbesserung aus
        Datenvielfalt, Datenmenge oder einer anderen Klassenverteilung
        resultiert; dafür sind kontrollierte Ablationen erforderlich.
\end{itemize}

\subsection{Grafikbedarf}
\begin{itemize}
  \item \textbf{Vorhanden:} \url{evaluation/plots/confusion_matrices_comparison_all_models.png}
        für den Hauptvergleich.
  \item \textbf{Vorhanden:} \url{evaluation/plots/roc_curves_comparison_all_models.png}
        zur Darstellung der ROC-Kurven.
  \item \textbf{Vorhanden:} \url{evaluation/plots/scenarios_all_confusion_matrices_grid.png}
        und die vorhandenen Szenariovergleiche im Ordner
        \url{evaluation/plots/}.
  \item \textbf{Erstellen:} Balkendiagramm für klassenweise Precision, Recall
        und F1; optional Falsch-Positiv-Rate nach deren Berechnung ergänzen.
\end{itemize}
